% GENERATED, edit the French unit or the translation table

\mxunit{stmg-c04}{Reciprocal (inverse) function}
\mxskills{finding a domain, differentiating, studying a sign, drawing up a table of variations, graphing}
\section{Reference notes}
The reciprocal function is defined by $f(x)=\dfrac1x$ on $\R\setminus\{0\}$.

\begin{center}
\begin{tikzpicture}
\begin{axis}[width=.75\textwidth,height=6.3cm,axis lines=middle,xmin=-5,xmax=5,ymin=-5,ymax=5,grid=major,xlabel={$x$},ylabel={$y$},unbounded coords=jump]
\addplot[mxteacher,very thick,samples=120,domain=-5:-.12]{1/x};
\addplot[mxteacher,very thick,samples=120,domain=.12:5]{1/x};
\addplot[mxwarn,dashed,thick] coordinates {(0,-5)(0,5)};
\addplot[mxwarn,dashed,thick] coordinates {(-5,0)(5,0)};
\end{axis}
\end{tikzpicture}
\end{center}

\begin{keypoint}
$f'(x)=-\dfrac1{x^2}<0$ for every $x\neq0$. The reciprocal function is strictly decreasing on each of the intervals $]-\infty,0[$ and $]0,+\infty[$. The lines $x=0$ and $y=0$ are asymptotes.
\end{keypoint}

\begin{pitfall}
One cannot write that the reciprocal function is decreasing ``on $\R\setminus\{0\}$'' without specifying the two intervals. For example $-1<1$ but $f(-1)=-1<1=f(1)$, which would contradict the global definition of decreasing.
\end{pitfall}

For a function $g(x)=P(x)+\dfrac{k}{x}$, proceed in this order: domain; derivative; common denominator and factorisation; sign of the derivative; variations; limits or behaviour at the boundaries; useful points; graph.

\begin{example}
On $]0,+\infty[$, $g(x)=x+\dfrac4x$ satisfies $g'(x)=1-\dfrac4{x^2}=\dfrac{(x-2)(x+2)}{x^2}$. The denominator is positive; $g'$ is negative on $]0,2[$ then positive on $]2,+\infty[$. The function has a minimum $g(2)=4$.
\end{example}

\mxsession{23}{\mxdate*{2026-12-15}}{B: 2 h -- C: 1 h}{Definition and graph of the reciprocal function}{Reciprocal (inverse) function}
\begin{objectives}\begin{itemize}\item Find the domain.\item Draw the curve.\item Use symmetries.\end{itemize}\end{objectives}
\begin{fluency}[6 to 8 minutes, no calculator]\begin{enumerate}\item Simplify $3x+5-x+2$. \mxblank{$2x+7$}\item Give the excluded value of $\dfrac{x+1}{x}$. \mxblank{$x=0$}\item Calculate $\dfrac1{5}$ and $\dfrac1{-5}$. \mxblank{$\dfrac1{5}$ and $-\dfrac1{5}$}\item Differentiate $3x^2-5x+1$. \mxblank{$6x-5$}\item Solve $(x-2)(x+1)=0$. \mxblank{$x=-1$ or $x=2$}\end{enumerate}\end{fluency}
\begin{prerequisites}Reread the last written record and be able to carry out the basic calculations needed today. Before the investigation, have students rephrase the key words: find the domain and draw the curve.\end{prerequisites}

\begin{activity}
Calculate $\dfrac1x$ for $x=-4,-2,-1,-\dfrac12,\dfrac12,1,2,4$. Plot the points and describe the symmetries. Why is $x=0$ missing?
\begin{center}\begin{tikzpicture}\begin{axis}[width=.68\textwidth,height=5cm,axis lines=middle,xmin=-4,xmax=4,ymin=-4,ymax=4,grid=major]\addplot[mxteacher,very thick,samples=80,domain=-4:-.15]{1/x};\addplot[mxteacher,very thick,samples=80,domain=.15:4]{1/x};\end{axis}\end{tikzpicture}\end{center}
\end{activity}
\begin{solution}The values are respectively $-\frac14,-\frac12,-1,-2,2,1,\frac12,\frac14$. The domain excludes $0$; the curve is symmetric about the origin. Have students say aloud which piece of data triggers the choice of method.\end{solution}

\begin{summary}The reciprocal function is $f(x)=\dfrac1x$ on $\R\setminus\{0\}$. Its curve is a hyperbola symmetric about the origin.\end{summary}

\begin{example}
Revisit the data from the investigation and write up the solution in order: definition of the objects, property used, exact or approximate calculation, then a concluding sentence. The model answer appears in the teacher edition.
\end{example}
\begin{solution}The values are respectively $-\frac14,-\frac12,-1,-2,2,1,\frac12,\frac14$. The domain excludes $0$; the curve is symmetric about the origin.\end{solution}

\begin{exercise}[essential]
\begin{enumerate}\item Give the domain and the asymptotes of $f(x)=5+\dfrac{7}x$.\item Calculate $f'(x)$ and study its sign.\item Study the function $g(x)=x+\dfrac{25}x$ on $]0,+\infty[$.\end{enumerate}
\end{exercise}
\begin{solution}The domain is $\R\setminus\{0\}$; asymptotes $x=0$ and $y=5$. $f'(x)=-\dfrac{7}{x^2}<0$. For $g$, $g'(x)=1-\dfrac{25}{x^2}$ and the minimum is reached at $x=5$, with value $10$.\end{solution}

\mxfurther{Calao, course pp. 88--93; exercises p. 94, no.~1--5}
\mxexuse{23}

\begin{homework}Redo practice exercise 2 with the following data: replace the initial value by $235$ and the rate or common ratio by $7$. Write it up in four lines: data, formula, calculation, conclusion.\end{homework}
\begin{solution}The answer depends on the exercise revisited. Check above all the consistency of the model, the number of iterations, the rounding and the concluding sentence.\end{solution}

\begin{exitticket}[3 minutes]On a half-sheet: write the main rule of the session, work through a very short example and state your confidence level from 1 to 4. Key question(s) of the day: 1, 2 and 3.\end{exitticket}

\ifmxteacher
\begin{method}
\textbf{1 h session:} 7 min of fluency; 10 min of pooling and written record; 28 min of essential exercises; 5 min of correction; 3 min of exit ticket.\\
\textbf{2 h session:} add 20 min of investigation, 20 min of formalisation, a second independent exercise and 10 min of comparative correction.\\
\textbf{Differentiation:} give struggling students a four-box sheet (data, model, calculation, conclusion). Ask quick students to produce a counterexample, a justification or a second method.
\end{method}
\fi
\mxsession{24}{\mxdate*{2026-12-17}}{B: 1 h -- C: 2 h}{Behaviour at the boundaries and asymptotes}{Reciprocal (inverse) function}
\begin{objectives}\begin{itemize}\item Describe the behaviour at the boundaries.\item Identify the asymptotes.\item Read a graph.\end{itemize}\end{objectives}
\begin{fluency}[6 to 8 minutes, no calculator]\begin{enumerate}\item Calculate $\pc{5}$ of $120$. \mxblank{$6$}\item Calculate $\dfrac1{6}$ and $\dfrac1{-6}$. \mxblank{$\dfrac1{6}$ and $-\dfrac1{6}$}\item Differentiate $3x^2-5x+1$. \mxblank{$6x-5$}\item Solve $(x-2)(x+1)=0$. \mxblank{$x=-1$ or $x=2$}\item Give the excluded value of $\dfrac{x+1}{x}$. \mxblank{$x=0$}\end{enumerate}\end{fluency}
\begin{prerequisites}Reread the last written record and be able to carry out the basic calculations needed today. Before the investigation, have students rephrase the key words: describe the behaviour at the boundaries and identify the asymptotes.\end{prerequisites}

\begin{activity}
Using a table of values near $0$ and for large values of $|x|$, describe $\dfrac1x$. Translate each behaviour into an asymptote and a sentence.
\begin{center}\begin{tikzpicture}\begin{axis}[width=.68\textwidth,height=5cm,axis lines=middle,xmin=-4,xmax=4,ymin=-4,ymax=4,grid=major]\addplot[mxteacher,very thick,samples=80,domain=-4:-.15]{1/x};\addplot[mxteacher,very thick,samples=80,domain=.15:4]{1/x};\end{axis}\end{tikzpicture}\end{center}
\end{activity}
\begin{solution}Near $0$, the values become very large in absolute value: asymptote $x=0$. As $|x|$ grows, $\dfrac1x$ gets closer to $0$: asymptote $y=0$. Have students say aloud which piece of data triggers the choice of method.\end{solution}

\begin{summary}For $f(x)=\dfrac1x$, $x=0$ is a vertical asymptote and $y=0$ a horizontal asymptote. The branches approach them without crossing them.\end{summary}

\begin{example}
Revisit the data from the investigation and write up the solution in order: definition of the objects, property used, exact or approximate calculation, then a concluding sentence. The model answer appears in the teacher edition.
\end{example}
\begin{solution}Near $0$, the values become very large in absolute value: asymptote $x=0$. As $|x|$ grows, $\dfrac1x$ gets closer to $0$: asymptote $y=0$.\end{solution}

\begin{exercise}[essential]
\begin{enumerate}\item Give the domain and the asymptotes of $f(x)=6+\dfrac{8}x$.\item Calculate $f'(x)$ and study its sign.\item Study the function $g(x)=x+\dfrac{36}x$ on $]0,+\infty[$.\end{enumerate}
\end{exercise}
\begin{solution}The domain is $\R\setminus\{0\}$; asymptotes $x=0$ and $y=6$. $f'(x)=-\dfrac{8}{x^2}<0$. For $g$, $g'(x)=1-\dfrac{36}{x^2}$ and the minimum is reached at $x=6$, with value $12$.\end{solution}

\mxfurther{Calao, course pp. 88--93; exercises p. 94, no.~4--9}
\mxexuse{24}

\begin{homework}Redo practice exercise 3 with the following data: replace the initial value by $240$ and the rate or common ratio by $2$. Write it up in four lines: data, formula, calculation, conclusion.\end{homework}
\begin{solution}The answer depends on the exercise revisited. Check above all the consistency of the model, the number of iterations, the rounding and the concluding sentence.\end{solution}

\begin{exitticket}[3 minutes]On a half-sheet: write the main rule of the session, work through a very short example and state your confidence level from 1 to 4. Key question(s) of the day: 1 and 2.\end{exitticket}

\ifmxteacher
\begin{method}
\textbf{1 h session:} 7 min of fluency; 10 min of pooling and written record; 28 min of essential exercises; 5 min of correction; 3 min of exit ticket.\\
\textbf{2 h session:} add 20 min of investigation, 20 min of formalisation, a second independent exercise and 10 min of comparative correction.\\
\textbf{Differentiation:} give struggling students a four-box sheet (data, model, calculation, conclusion). Ask quick students to produce a counterexample, a justification or a second method.
\end{method}
\fi
\mxsession{25}{\mxdate*{2027-01-05}}{B: 2 h -- C: 1 h}{Derivative and direction of variation}{Reciprocal (inverse) function}
\begin{objectives}\begin{itemize}\item Differentiate the reciprocal function.\item Justify the variations.\item Draw up a table.\end{itemize}\end{objectives}
\begin{fluency}[6 to 8 minutes, no calculator]\begin{enumerate}\item Write the multiplier for an increase of $\pc{10}$. \mxblank{$1+\dfrac{10}{100}=\num{1,10}$}\item Differentiate $3x^2-5x+1$. \mxblank{$6x-5$}\item Solve $(x-2)(x+1)=0$. \mxblank{$x=-1$ or $x=2$}\item Give the excluded value of $\dfrac{x+1}{x}$. \mxblank{$x=0$}\item Calculate $\dfrac1{7}$ and $\dfrac1{-7}$. \mxblank{$\dfrac1{7}$ and $-\dfrac1{7}$}\end{enumerate}\end{fluency}
\begin{prerequisites}Reread the last written record and be able to carry out the basic calculations needed today. Before the investigation, have students rephrase the key words: differentiate the reciprocal function and justify the variations.\end{prerequisites}

\begin{activity}
Starting from the difference quotient of $f(x)=\dfrac1x$, recover $f'(x)=-\dfrac1{x^2}$. Study its sign and check consistency with the graph.

\end{activity}
\begin{solution}The square $x^2$ is positive for $x\neq0$, so the derivative is strictly negative on each interval of the domain. Have students say aloud which piece of data triggers the choice of method.\end{solution}

\begin{summary}For $x\neq0$, $f'(x)=-\dfrac1{x^2}<0$. The reciprocal function decreases on each of the intervals $]-\infty,0[$ and $]0,+\infty[$.\end{summary}

\begin{example}
Revisit the data from the investigation and write up the solution in order: definition of the objects, property used, exact or approximate calculation, then a concluding sentence. The model answer appears in the teacher edition.
\end{example}
\begin{solution}The square $x^2$ is positive for $x\neq0$, so the derivative is strictly negative on each interval of the domain.\end{solution}

\begin{exercise}[essential]
\begin{enumerate}\item Give the domain and the asymptotes of $f(x)=2+\dfrac{4}x$.\item Calculate $f'(x)$ and study its sign.\item Study the function $g(x)=x+\dfrac{4}x$ on $]0,+\infty[$.\end{enumerate}
\end{exercise}
\begin{solution}The domain is $\R\setminus\{0\}$; asymptotes $x=0$ and $y=2$. $f'(x)=-\dfrac{4}{x^2}<0$. For $g$, $g'(x)=1-\dfrac{4}{x^2}$ and the minimum is reached at $x=2$, with value $4$.\end{solution}

\mxfurther{Calao, course pp. 89--93; exercises p. 94, no.~1--3, then p. 96, no.~10--13}
\mxexuse{25}

\begin{homework}Redo practice exercise 2 with the following data: replace the initial value by $245$ and the rate or common ratio by $3$. Write it up in four lines: data, formula, calculation, conclusion.\end{homework}
\begin{solution}The answer depends on the exercise revisited. Check above all the consistency of the model, the number of iterations, the rounding and the concluding sentence.\end{solution}

\begin{exitticket}[3 minutes]On a half-sheet: write the main rule of the session, work through a very short example and state your confidence level from 1 to 4. Key question(s) of the day: 1, 2 and 3.\end{exitticket}

\ifmxteacher
\begin{method}
\textbf{1 h session:} 7 min of fluency; 10 min of pooling and written record; 28 min of essential exercises; 5 min of correction; 3 min of exit ticket.\\
\textbf{2 h session:} add 20 min of investigation, 20 min of formalisation, a second independent exercise and 10 min of comparative correction.\\
\textbf{Differentiation:} give struggling students a four-box sheet (data, model, calculation, conclusion). Ask quick students to produce a counterexample, a justification or a second method.
\end{method}
\fi
\mxsession{26}{\mxdate*{2027-01-07}}{B: 1 h -- C: 2 h}{Combinations with polynomial functions}{Reciprocal (inverse) function}
\begin{objectives}\begin{itemize}\item Differentiate a combination of functions.\item Simplify an expression.\item Organise a function study.\end{itemize}\end{objectives}
\begin{fluency}[6 to 8 minutes, no calculator]\begin{enumerate}\item Simplify $3x+8-x+2$. \mxblank{$2x+10$}\item Solve $(x-2)(x+1)=0$. \mxblank{$x=-1$ or $x=2$}\item Give the excluded value of $\dfrac{x+1}{x}$. \mxblank{$x=0$}\item Calculate $\dfrac1{8}$ and $\dfrac1{-8}$. \mxblank{$\dfrac1{8}$ and $-\dfrac1{8}$}\item Differentiate $3x^2-5x+1$. \mxblank{$6x-5$}\end{enumerate}\end{fluency}
\begin{prerequisites}Reread the last written record and be able to carry out the basic calculations needed today. Before the investigation, have students rephrase the key words: differentiate a combination of functions and simplify an expression.\end{prerequisites}

\begin{activity}
A company produces $x>0$ hundred items and its average cost is $C(x)=2x+30+\dfrac{200}{x}$. Calculate $C'(x)$, factorise its numerator and find the production level that minimises the average cost.

\end{activity}
\begin{solution}$C'(x)=2-\dfrac{200}{x^2}=\dfrac{2(x-10)(x+10)}{x^2}$. The minimum is reached for $x=10$, that is $\num{1000}$ items. Have students say aloud which piece of data triggers the choice of method.\end{solution}

\begin{summary}To study $P(x)+\dfrac{k}{x}$, find the domain, calculate the derivative, put over a common denominator, study its sign, then draw up the variations.\end{summary}

\begin{example}
Revisit the data from the investigation and write up the solution in order: definition of the objects, property used, exact or approximate calculation, then a concluding sentence. The model answer appears in the teacher edition.
\end{example}
\begin{solution}$C'(x)=2-\dfrac{200}{x^2}=\dfrac{2(x-10)(x+10)}{x^2}$. The minimum is reached for $x=10$, that is $\num{1000}$ items.\end{solution}

\begin{exercise}[essential]
\begin{enumerate}\item Give the domain and the asymptotes of $f(x)=3+\dfrac{5}x$.\item Calculate $f'(x)$ and study its sign.\item Study the function $g(x)=x+\dfrac{9}x$ on $]0,+\infty[$.\end{enumerate}
\end{exercise}
\begin{solution}The domain is $\R\setminus\{0\}$; asymptotes $x=0$ and $y=3$. $f'(x)=-\dfrac{5}{x^2}<0$. For $g$, $g'(x)=1-\dfrac{9}{x^2}$ and the minimum is reached at $x=3$, with value $6$.\end{solution}

\mxfurther{Calao, p. 96, no.~14--21}
\mxexuse{26}

\begin{homework}Redo practice exercise 3 with the following data: replace the initial value by $250$ and the rate or common ratio by $4$. Write it up in four lines: data, formula, calculation, conclusion.\end{homework}
\begin{solution}The answer depends on the exercise revisited. Check above all the consistency of the model, the number of iterations, the rounding and the concluding sentence.\end{solution}

\begin{exitticket}[3 minutes]On a half-sheet: write the main rule of the session, work through a very short example and state your confidence level from 1 to 4. Key question(s) of the day: 1 and 2.\end{exitticket}

\ifmxteacher
\begin{method}
\textbf{1 h session:} 7 min of fluency; 10 min of pooling and written record; 28 min of essential exercises; 5 min of correction; 3 min of exit ticket.\\
\textbf{2 h session:} add 20 min of investigation, 20 min of formalisation, a second independent exercise and 10 min of comparative correction.\\
\textbf{Differentiation:} give struggling students a four-box sheet (data, model, calculation, conclusion). Ask quick students to produce a counterexample, a justification or a second method.
\end{method}
\fi
\mxsession{27}{\mxdate*{2027-01-12}}{B: 2 h -- C: 1 h}{Sign and factorisation}{Reciprocal (inverse) function}
\begin{objectives}\begin{itemize}\item Factorise.\item Build a sign table.\item Respect excluded values.\end{itemize}\end{objectives}
\begin{fluency}[6 to 8 minutes, no calculator]\begin{enumerate}\item Calculate $\pc{20}$ of $180$. \mxblank{$36$}\item Give the excluded value of $\dfrac{x+1}{x}$. \mxblank{$x=0$}\item Calculate $\dfrac1{9}$ and $\dfrac1{-9}$. \mxblank{$\dfrac1{9}$ and $-\dfrac1{9}$}\item Differentiate $3x^2-5x+1$. \mxblank{$6x-5$}\item Solve $(x-2)(x+1)=0$. \mxblank{$x=-1$ or $x=2$}\end{enumerate}\end{fluency}
\begin{prerequisites}Reread the last written record and be able to carry out the basic calculations needed today. Before the investigation, have students rephrase the key words: factorise and build a sign table.\end{prerequisites}

\begin{activity}
Study the sign of $\dfrac{(x-2)(x+3)}x$. Build the table with the three critical values, distinguish a zero from an excluded value, then solve the inequality $\dfrac{(x-2)(x+3)}x\geq0$.

\end{activity}
\begin{solution}The critical values are $-3,0,2$. The solution is $[-3,0[\cup[2,+\infty[$. Have students say aloud which piece of data triggers the choice of method.\end{solution}

\begin{summary}In a quotient, the sign depends on the numerator and the denominator. Excluded values are left out of the table, even if a factor formally cancels.\end{summary}

\begin{example}
Revisit the data from the investigation and write up the solution in order: definition of the objects, property used, exact or approximate calculation, then a concluding sentence. The model answer appears in the teacher edition.
\end{example}
\begin{solution}The critical values are $-3,0,2$. The solution is $[-3,0[\cup[2,+\infty[$.\end{solution}

\begin{exercise}[essential]
\begin{enumerate}\item Give the domain and the asymptotes of $f(x)=4+\dfrac{6}x$.\item Calculate $f'(x)$ and study its sign.\item Study the function $g(x)=x+\dfrac{16}x$ on $]0,+\infty[$.\end{enumerate}
\end{exercise}
\begin{solution}The domain is $\R\setminus\{0\}$; asymptotes $x=0$ and $y=4$. $f'(x)=-\dfrac{6}{x^2}<0$. For $g$, $g'(x)=1-\dfrac{16}{x^2}$ and the minimum is reached at $x=4$, with value $8$.\end{solution}

\mxfurther{Calao, pp. 96--101}
\mxexuse{27}

\begin{homework}Redo practice exercise 2 with the following data: replace the initial value by $255$ and the rate or common ratio by $5$. Write it up in four lines: data, formula, calculation, conclusion.\end{homework}
\begin{solution}The answer depends on the exercise revisited. Check above all the consistency of the model, the number of iterations, the rounding and the concluding sentence.\end{solution}

\begin{exitticket}[3 minutes]On a half-sheet: write the main rule of the session, work through a very short example and state your confidence level from 1 to 4. Key question(s) of the day: 1, 2 and 3.\end{exitticket}

\ifmxteacher
\begin{method}
\textbf{1 h session:} 7 min of fluency; 10 min of pooling and written record; 28 min of essential exercises; 5 min of correction; 3 min of exit ticket.\\
\textbf{2 h session:} add 20 min of investigation, 20 min of formalisation, a second independent exercise and 10 min of comparative correction.\\
\textbf{Differentiation:} give struggling students a four-box sheet (data, model, calculation, conclusion). Ask quick students to produce a counterexample, a justification or a second method.
\end{method}
\fi
\mxsession{28}{\mxdate*{2027-01-14}}{B: 1 h -- C: 2 h}{Derivatives of polynomials of degree at most 3}{Reciprocal (inverse) function}
\begin{objectives}\begin{itemize}\item Differentiate a polynomial.\item Study a derivative.\item Deduce the variations.\end{itemize}\end{objectives}
\begin{fluency}[6 to 8 minutes, no calculator]\begin{enumerate}\item Write the multiplier for an increase of $\pc{25}$. \mxblank{$1+\dfrac{25}{100}=\num{1,25}$}\item Calculate $\dfrac1{3}$ and $\dfrac1{-3}$. \mxblank{$\dfrac1{3}$ and $-\dfrac1{3}$}\item Differentiate $3x^2-5x+1$. \mxblank{$6x-5$}\item Solve $(x-2)(x+1)=0$. \mxblank{$x=-1$ or $x=2$}\item Give the excluded value of $\dfrac{x+1}{x}$. \mxblank{$x=0$}\end{enumerate}\end{fluency}
\begin{prerequisites}Reread the last written record and be able to carry out the basic calculations needed today. Before the investigation, have students rephrase the key words: differentiate a polynomial and study a derivative.\end{prerequisites}

\begin{activity}
For $f(x)=2x^3-3x^2-12x+5$, calculate $f'(x)$, factorise, study its sign and deduce the variations of $f$.

\end{activity}
\begin{solution}$f'(x)=6x^2-6x-12=6(x-2)(x+1)$. The function increases on $]-\infty,-1]$, decreases on $[-1,2]$, then increases. Have students say aloud which piece of data triggers the choice of method.\end{solution}

\begin{summary}$(ax^3+bx^2+cx+d)'=3ax^2+2bx+c$. The sign of the derivative gives the direction of variation of the function.\end{summary}

\begin{example}
Revisit the data from the investigation and write up the solution in order: definition of the objects, property used, exact or approximate calculation, then a concluding sentence. The model answer appears in the teacher edition.
\end{example}
\begin{solution}$f'(x)=6x^2-6x-12=6(x-2)(x+1)$. The function increases on $]-\infty,-1]$, decreases on $[-1,2]$, then increases.\end{solution}

\begin{exercise}[essential]
\begin{enumerate}\item Give the domain and the asymptotes of $f(x)=5+\dfrac{7}x$.\item Calculate $f'(x)$ and study its sign.\item Study the function $g(x)=x+\dfrac{25}x$ on $]0,+\infty[$.\end{enumerate}
\end{exercise}
\begin{solution}The domain is $\R\setminus\{0\}$; asymptotes $x=0$ and $y=5$. $f'(x)=-\dfrac{7}{x^2}<0$. For $g$, $g'(x)=1-\dfrac{25}{x^2}$ and the minimum is reached at $x=5$, with value $10$.\end{solution}

\mxfurther{Calao, pp. 94--96}
\mxexuse{28}

\begin{homework}Redo practice exercise 3 with the following data: replace the initial value by $260$ and the rate or common ratio by $6$. Write it up in four lines: data, formula, calculation, conclusion.\end{homework}
\begin{solution}The answer depends on the exercise revisited. Check above all the consistency of the model, the number of iterations, the rounding and the concluding sentence.\end{solution}

\begin{exitticket}[3 minutes]On a half-sheet: write the main rule of the session, work through a very short example and state your confidence level from 1 to 4. Key question(s) of the day: 1 and 2.\end{exitticket}

\ifmxteacher
\begin{method}
\textbf{1 h session:} 7 min of fluency; 10 min of pooling and written record; 28 min of essential exercises; 5 min of correction; 3 min of exit ticket.\\
\textbf{2 h session:} add 20 min of investigation, 20 min of formalisation, a second independent exercise and 10 min of comparative correction.\\
\textbf{Differentiation:} give struggling students a four-box sheet (data, model, calculation, conclusion). Ask quick students to produce a counterexample, a justification or a second method.
\end{method}
\fi
\mxsession{29}{\mxdate*{2027-01-19}}{B: 2 h -- C: 1 h}{Full study and graph}{Reciprocal (inverse) function}
\begin{objectives}\begin{itemize}\item Carry out a full study.\item Draw a consistent curve.\item Interpret an extremum.\end{itemize}\end{objectives}
\begin{fluency}[6 to 8 minutes, no calculator]\begin{enumerate}\item Simplify $3x+4-x+2$. \mxblank{$2x+6$}\item Differentiate $3x^2-5x+1$. \mxblank{$6x-5$}\item Solve $(x-2)(x+1)=0$. \mxblank{$x=-1$ or $x=2$}\item Give the excluded value of $\dfrac{x+1}{x}$. \mxblank{$x=0$}\item Calculate $\dfrac1{4}$ and $\dfrac1{-4}$. \mxblank{$\dfrac1{4}$ and $-\dfrac1{4}$}\end{enumerate}\end{fluency}
\begin{prerequisites}Reread the last written record and be able to carry out the basic calculations needed today. Before the investigation, have students rephrase the key words: carry out a full study and draw a consistent curve.\end{prerequisites}

\begin{activity}
Study $g(x)=x+\dfrac9x$ completely on $]0,+\infty[$: domain, derivative, sign, variations, minimum, behaviour at the boundaries and sketch.
\begin{center}\begin{tikzpicture}\begin{axis}[width=.68\textwidth,height=5cm,axis lines=middle,xmin=-4,xmax=4,ymin=-4,ymax=4,grid=major]\addplot[mxteacher,very thick,samples=80,domain=-4:-.15]{1/x};\addplot[mxteacher,very thick,samples=80,domain=.15:4]{1/x};\end{axis}\end{tikzpicture}\end{center}
\end{activity}
\begin{solution}$g'(x)=\dfrac{(x-3)(x+3)}{x^2}$; minimum $g(3)=6$; $g(x)$ tends to $+\infty$ at $0^+$ and at $+\infty$. Have students say aloud which piece of data triggers the choice of method.\end{solution}

\begin{summary}A full study includes: domain, derivative, sign, variations, limits or behaviour at the boundaries, key points and a consistent graph.\end{summary}

\begin{example}
Revisit the data from the investigation and write up the solution in order: definition of the objects, property used, exact or approximate calculation, then a concluding sentence. The model answer appears in the teacher edition.
\end{example}
\begin{solution}$g'(x)=\dfrac{(x-3)(x+3)}{x^2}$; minimum $g(3)=6$; $g(x)$ tends to $+\infty$ at $0^+$ and at $+\infty$.\end{solution}

\begin{exercise}[essential]
\begin{enumerate}\item Give the domain and the asymptotes of $f(x)=6+\dfrac{8}x$.\item Calculate $f'(x)$ and study its sign.\item Study the function $g(x)=x+\dfrac{36}x$ on $]0,+\infty[$.\end{enumerate}
\end{exercise}
\begin{solution}The domain is $\R\setminus\{0\}$; asymptotes $x=0$ and $y=6$. $f'(x)=-\dfrac{8}{x^2}<0$. For $g$, $g'(x)=1-\dfrac{36}{x^2}$ and the minimum is reached at $x=6$, with value $12$.\end{solution}

\mxfurther{Calao, pp. 98--103}
\mxexuse{29}

\begin{homework}Redo practice exercise 2 with the following data: replace the initial value by $265$ and the rate or common ratio by $7$. Write it up in four lines: data, formula, calculation, conclusion.\end{homework}
\begin{solution}The answer depends on the exercise revisited. Check above all the consistency of the model, the number of iterations, the rounding and the concluding sentence.\end{solution}

\begin{exitticket}[3 minutes]On a half-sheet: write the main rule of the session, work through a very short example and state your confidence level from 1 to 4. Key question(s) of the day: 1, 2 and 3.\end{exitticket}

\ifmxteacher
\begin{method}
\textbf{1 h session:} 7 min of fluency; 10 min of pooling and written record; 28 min of essential exercises; 5 min of correction; 3 min of exit ticket.\\
\textbf{2 h session:} add 20 min of investigation, 20 min of formalisation, a second independent exercise and 10 min of comparative correction.\\
\textbf{Differentiation:} give struggling students a four-box sheet (data, model, calculation, conclusion). Ask quick students to produce a counterexample, a justification or a second method.
\end{method}
\fi
\mxsession{30}{\mxdate*{2027-01-21}}{B: 1 h -- C: 2 h}{Review and assessment}{Reciprocal (inverse) function}
\begin{objectives}\begin{itemize}\item Use domain, derivative and sign.\item Write up a study.\item Correct work methodically.\end{itemize}\end{objectives}
\begin{fluency}[6 to 8 minutes, no calculator]\begin{enumerate}\item Calculate $\pc{5}$ of $100$. \mxblank{$5$}\item Solve $(x-2)(x+1)=0$. \mxblank{$x=-1$ or $x=2$}\item Give the excluded value of $\dfrac{x+1}{x}$. \mxblank{$x=0$}\item Calculate $\dfrac1{5}$ and $\dfrac1{-5}$. \mxblank{$\dfrac1{5}$ and $-\dfrac1{5}$}\item Differentiate $3x^2-5x+1$. \mxblank{$6x-5$}\end{enumerate}\end{fluency}
\begin{prerequisites}Reread the last written record and be able to carry out the basic calculations needed today. Before the investigation, have students rephrase the key words: use domain, derivative and sign and write up a study.\end{prerequisites}

\begin{activity}
Correct a deliberately flawed study of $h(x)=x+\dfrac4x$: stated domain $\R$, derivative $1+\dfrac4{x^2}$, table with no break at $0$. Rewrite a complete solution.

\end{activity}
\begin{solution}The domain is $\R\setminus\{0\}$ and $h'(x)=1-\dfrac4{x^2}$. The table is split at $0$ and takes into account $-2$ and $2$. Have students say aloud which piece of data triggers the choice of method.\end{solution}

\begin{summary}Check the four items separately: domain, expression of the derivative, sign of the derivative, variations. An error at one step must not hide the following ones.\end{summary}

\begin{example}
Revisit the data from the investigation and write up the solution in order: definition of the objects, property used, exact or approximate calculation, then a concluding sentence. The model answer appears in the teacher edition.
\end{example}
\begin{solution}The domain is $\R\setminus\{0\}$ and $h'(x)=1-\dfrac4{x^2}$. The table is split at $0$ and takes into account $-2$ and $2$.\end{solution}

\begin{exercise}[essential]
\begin{enumerate}\item Give the domain and the asymptotes of $f(x)=2+\dfrac{4}x$.\item Calculate $f'(x)$ and study its sign.\item Study the function $g(x)=x+\dfrac{4}x$ on $]0,+\infty[$.\end{enumerate}
\end{exercise}
\begin{solution}The domain is $\R\setminus\{0\}$; asymptotes $x=0$ and $y=2$. $f'(x)=-\dfrac{4}{x^2}<0$. For $g$, $g'(x)=1-\dfrac{4}{x^2}$ and the minimum is reached at $x=2$, with value $4$.\end{solution}

\mxfurther{Calao, pp. 100--103}
\mxexuse{30}

\begin{homework}Redo practice exercise 3 with the following data: replace the initial value by $270$ and the rate or common ratio by $2$. Write it up in four lines: data, formula, calculation, conclusion.\end{homework}
\begin{solution}The answer depends on the exercise revisited. Check above all the consistency of the model, the number of iterations, the rounding and the concluding sentence.\end{solution}

\begin{exitticket}[3 minutes]On a half-sheet: write the main rule of the session, work through a very short example and state your confidence level from 1 to 4. Key question(s) of the day: 1 and 2.\end{exitticket}

\ifmxteacher
\begin{method}
\textbf{1 h session:} 7 min of fluency; 10 min of pooling and written record; 28 min of essential exercises; 5 min of correction; 3 min of exit ticket.\\
\textbf{2 h session:} add 20 min of investigation, 20 min of formalisation, a second independent exercise and 10 min of comparative correction.\\
\textbf{Differentiation:} give struggling students a four-box sheet (data, model, calculation, conclusion). Ask quick students to produce a counterexample, a justification or a second method.
\end{method}
\fi
